%% ============================================================
%%  5.  EXPERIMENTAL SETUP
%% ============================================================
\section{Experimental Setup}
\label{sec:experiments}

\paragraph{Dataset.}
We use the \textbf{MMLU-Pro} benchmark~\cite{mmlupro}, a challenging
multiple-choice reasoning dataset spanning 14 academic domains (biology,
chemistry, law, mathematics, etc.) with up to 10 answer options per
question.
We sample 200 questions from the test split at a fixed seed (seed=7)
and run each experimental condition on the same 200 questions, enabling
paired statistical comparison.
All questions are in categorical (multiple-choice) mode; no ordinal stance
scores are used.

\paragraph{Model.}
Experiments 1--3b use \textbf{Qwen2.5-14B-Instruct}~\cite{qwen25} via the
HuggingFace \texttt{transformers} backend on a cluster GPU.
Experiment~4 uses \textbf{Qwen2.5-7B-Instruct} via Ollama for local
reproducibility on consumer hardware (NVIDIA RTX A2000 8\,GB).
All agents use temperature $0.7$, top-$p$ $0.9$, and a maximum of 220
new tokens per turn (extended for protected agents; see \S\ref{sec:predictor}).

\paragraph{Debate protocol.}
Each question is debated by $N\!=\!3$ agents over $T\!=\!5$ rounds.
Round~1 agents answer independently.
Rounds 2--5 agents see all previous answers and explanations before
responding.
If all agents agree by round $T$ the consensus answer is accepted; otherwise
a moderator selects the majority answer.
Agent presentation order is shuffled per question.
Parse failures trigger up to two re-prompts; debates with unresolved parse
failures are discarded.

\paragraph{LLM-as-judge.}
After all debates complete, the same Qwen model evaluates each agent
explanation along four dimensions (explanation quality, use of prior-round
reasoning, justification of stance, independence of reasoning), each scored
in $[0,1]$.
The overall score $q_i^{(t)}$ used in feature engineering and process
metrics is the \texttt{explanation\_good} dimension.

\paragraph{Evaluation metrics.}
We report:
\textbf{Accuracy} (fraction of correct final answers),
\textbf{PER} (Eq.~2),
\textbf{Consensus Reliability} ($P(\text{correct}\mid\text{consensus})$),
and the six process-level diagnostics from \citeauthor{diagnostic2026}:
engagement, responsiveness, influence asymmetry, balance, stability, and
group welfare (all computed with categorical adaptations).

\paragraph{Conditions.}
\begin{itemize}[noitemsep,topsep=2pt]
  \item \textbf{Baseline} (Exp1): standard debate, no intervention.
  \item \textbf{Loser Resistance--Support} (Exp2a): support-trend trigger.
  \item \textbf{Loser Resistance--Metric} (Exp2b): diagnostic trigger.
  \item \textbf{Static Minority Protection} (Exp3a): categorical trigger.
  \item \textbf{Diagnostic Protection} (Exp3b): diagnostic trigger (\S\ref{sec:mechanisms}).
  \item \textbf{Learned Predictor} (Exp4): calibrated logistic predictor (\S\ref{sec:predictor}).
\end{itemize}
Exp3b and Exp4 share identical prompt instructions for flagged agents; the
only difference is how the flag is set.
This isolates the contribution of the learned predictor over the hand-tuned
thresholds.

\paragraph{Statistical tests.}
We use McNemar's test for paired accuracy comparisons (same 200 questions)
and Wilcoxon signed-rank tests for process metrics.
Significance levels: $^{*}p\!<\!0.05$, $^{**}p\!<\!0.01$,
$^{***}p\!<\!0.001$.
